\section*{Lecture 1: Overview}
\addcontentsline{toc}{section}{Lecture 1: Overview}
\clearpage
\SlideHeader{Lecture 1: Overview}{001}{表紙}
\SlideImage{1}{../Materials/2026/3DM_01_Overview_SS26.pdf}
\begin{multicols}{2}\scriptsize\raggedcolumns
\NoteSection{日本語訳}
\begin{itemize}
\item Data Driven Decision Making
\item 第1回 - 概要
\end{itemize}
\end{multicols}

\clearpage
\SlideHeader{Lecture 1: Overview}{002}{連絡先}
\SlideImage{2}{../Materials/2026/3DM_01_Overview_SS26.pdf}
\begin{multicols}{2}\scriptsize\raggedcolumns
\NoteSection{日本語訳}
\begin{itemize}
\item 担当教員
\begin{itemize}
\item Prof. Dr. Dirk Mattfeld
\item オフィスアワー: 登録制
\item d.mattfeld@tu-braunschweig.de
\end{itemize}
\item ティーチングアシスタント
\begin{itemize}
\item Felix Moeller
\item オフィスアワー: 登録制
\item felix.moeller@tu-braunschweig.de
\end{itemize}
\item Decision Support オフィス
\begin{itemize}
\item オフィスアワー: Webサイトで告知
\item 電話: (0531) 391-3211
\item ds@tu-braunschweig.de
\end{itemize}
\item Webサイト
\begin{itemize}
\item 講義・試験日程などを確認する場所
\end{itemize}
\end{itemize}
\end{multicols}

\clearpage
\SlideHeader{Lecture 1: Overview}{003}{Decision Support の修士科目}
\SlideImage{3}{../Materials/2026/3DM_01_Overview_SS26.pdf}
\begin{multicols}{2}\scriptsize\raggedcolumns
\NoteSection{日本語訳}
\begin{itemize}
\item 基礎・方向づけ
\begin{itemize}
\item Decision Support (5 LP)
\item Intelligent Data Analysis (WS)
\item Planning for Mobility and Transportation Purposes (WS)
\end{itemize}
\item 専門
\begin{itemize}
\item Decision Support (5 LP)
\item Data Driven Decision Making (SS)
\item Exercise in Data Driven Decision Making (SS)
\item Master seminar
\item Master thesis (6 months)
\end{itemize}
\item その他の専攻、とくに機械工学系学生にも開かれている。
\end{itemize}
\end{multicols}

\clearpage
\SlideHeader{Lecture 1: Overview}{004}{講義について}
\SlideImage{4}{../Materials/2026/3DM_01_Overview_SS26.pdf}
\begin{multicols}{2}\footnotesize\raggedcolumns
\NoteSection{日本語訳}
\begin{itemize}
\item Data Driven Decision Making
\item 以前は『モビリティ応用のための情報システム』、さらに以前は『交通情報システム』だった。
\item 焦点は、情報システムの描写から、情報システム内の自動意思決定へ移っている。
\item Uber, Lime, Amazon などの先端的デジタルビジネスモデルと整合する。
\item PMT は意思決定に、IDA はデータ分析に焦点を置く。DDDM はその上に構築される。
\item 3DM は演習とともに『Decision Support専門』を構成する。
\end{itemize}
\NoteSection{解説}
DDDMでは、データを眺めるだけでなく、データに基づいて実行可能な行動を選ぶ。例えば、食品配送なら『どの注文をどのドライバーに割り当てるか』、在庫管理なら『いつ何個発注するか』が決定である。\par

stochasticity は、将来情報が確実には分からないことを意味する。例として、注文発生、道路混雑、サービス時間、ドライバーのキャンセルがある。dynamism は、新しい情報が入るたびに状態と計画が変わることである。\par

試験答案では、stochasticity=不確実な情報実現、dynamism=時間を通じた状態変化と再意思決定、と分けて書き、最後に『そのため将来を先読みしつつ柔軟性を残す必要がある』と結ぶとよい。\par
\NoteSection{試験対策ポイント}
\begin{itemize}
\item DDDM = データ分析 + 意思決定モデル + 現実への実装、というループで理解する。
\item 単なる予測ではなく、予測・観測を意思決定に変換する点が講義の中心である。
\end{itemize}
\NoteSection{確認問題と解答}
\begin{enumerate}
\item \textbf{DDDMは通常のデータ分析と何が違うか。}\\通常のデータ分析は、過去データの記述や将来予測で終わる場合がある。DDDMでは、その予測や観測を使って具体的な行動を選ぶ。配送なら需要予測だけでなく、車両割当、訪問順序、注文受諾、再計画まで含める。したがって情報モデル、意思決定モデル、実装を一体で説明する必要がある。
\item \textbf{stochasticity と dynamism を例で区別せよ。}\\食品配送で『次にどこから注文が来るか分からない』ことは stochasticity である。注文が実際に入った後、ドライバー位置や未配送注文が変わり、計画を作り直す必要があることは dynamism である。両者が合わさると stochastic dynamic decision problem になる。
\end{enumerate}
\end{multicols}

\clearpage
\SlideHeader{Lecture 1: Overview}{005}{目標}
\SlideImage{5}{../Materials/2026/3DM_01_Overview_SS26.pdf}
\begin{multicols}{2}\footnotesize\raggedcolumns
\NoteSection{日本語訳}
\begin{itemize}
\item 確率性と動的性質を理解する。
\item 確率的・動的意思決定問題をモデル化する。
\item 方法論の概要、機能、強み、弱みを理解する。
\item デジタル経済における関連問題を理解する。
\item 60分筆記試験に合格する。
\item 修士論文、博士課程への準備を行う。
\end{itemize}
\NoteSection{解説}
DDDMでは、データを眺めるだけでなく、データに基づいて実行可能な行動を選ぶ。例えば、食品配送なら『どの注文をどのドライバーに割り当てるか』、在庫管理なら『いつ何個発注するか』が決定である。\par

stochasticity は、将来情報が確実には分からないことを意味する。例として、注文発生、道路混雑、サービス時間、ドライバーのキャンセルがある。dynamism は、新しい情報が入るたびに状態と計画が変わることである。\par

試験答案では、stochasticity=不確実な情報実現、dynamism=時間を通じた状態変化と再意思決定、と分けて書き、最後に『そのため将来を先読みしつつ柔軟性を残す必要がある』と結ぶとよい。\par
\NoteSection{試験対策ポイント}
\begin{itemize}
\item DDDM = データ分析 + 意思決定モデル + 現実への実装、というループで理解する。
\item 単なる予測ではなく、予測・観測を意思決定に変換する点が講義の中心である。
\end{itemize}
\NoteSection{確認問題と解答}
\begin{enumerate}
\item \textbf{DDDMは通常のデータ分析と何が違うか。}\\通常のデータ分析は、過去データの記述や将来予測で終わる場合がある。DDDMでは、その予測や観測を使って具体的な行動を選ぶ。配送なら需要予測だけでなく、車両割当、訪問順序、注文受諾、再計画まで含める。したがって情報モデル、意思決定モデル、実装を一体で説明する必要がある。
\item \textbf{stochasticity と dynamism を例で区別せよ。}\\食品配送で『次にどこから注文が来るか分からない』ことは stochasticity である。注文が実際に入った後、ドライバー位置や未配送注文が変わり、計画を作り直す必要があることは dynamism である。両者が合わさると stochastic dynamic decision problem になる。
\end{enumerate}
\end{multicols}

\clearpage
\SlideHeader{Lecture 1: Overview}{006}{DDDMの講義計画}
\SlideImage{6}{../Materials/2026/3DM_01_Overview_SS26.pdf}
\begin{multicols}{2}\scriptsize\raggedcolumns
\NoteSection{日本語訳}
\begin{itemize}
\item 講義トピック
\begin{itemize}
\item 1 概要
\item 2 情報モデリングと意思決定モデリング
\item 3 確率性
\item 4 動的性
\item 5 マルコフ意思決定過程
\item 6 近似動的計画法
\item 7 Look-ahead Methods
\item 8 価値関数近似
\item 9 ADPにおける分析上の問題
\item 10 Combined Methods
\item 11 プラットフォーム型フードデリバリー
\end{itemize}
\end{itemize}
\end{multicols}

\clearpage
\SlideHeader{Lecture 1: Overview}{007}{文献}
\SlideImage{7}{../Materials/2026/3DM_01_Overview_SS26.pdf}
\begin{multicols}{2}\scriptsize\raggedcolumns
\NoteSection{日本語訳}
\begin{itemize}
\item Warren Powell の研究。
\item TU Braunschweig における Marlin Ulmer らとの近年の共同研究。
\item 交通・モビリティに関する適用と発展。
\item Ulmer et al. (2020): stochastic dynamic vehicle routing problems のモデル化。
\item Ulmer et al. (2021): prescriptive analytics の観点から見た stochastic dynamic vehicle routing のレビュー。
\end{itemize}
\end{multicols}

\clearpage
\SlideHeader{Lecture 1: Overview}{008}{演習について}
\SlideImage{8}{../Materials/2026/3DM_01_Overview_SS26.pdf}
\begin{multicols}{2}\scriptsize\raggedcolumns
\NoteSection{日本語訳}
\begin{itemize}
\item 『Decision Support専門』の Studienleistung (2.5 LP)。
\item serious game アプリケーション『Trucks and Barges』で講義内容を練習する。
\item シミュレーション内容
\begin{itemize}
\item 貨物業者が、複数日にわたり、ターミナルからコンテナを異なる輸送モードで出荷する意思決定状況。
\end{itemize}
\item 運営
\begin{itemize}
\item 内容が積み上がる5回のオンライン日程。
\item 宿題。
\item StudIP: Data Driven Decision Making - Uebung
\end{itemize}
\end{itemize}
\end{multicols}

\clearpage
\SlideHeader{Lecture 1: Overview}{009}{試験情報}
\SlideImage{9}{../Materials/2026/3DM_01_Overview_SS26.pdf}
\begin{multicols}{2}\scriptsize\raggedcolumns
\NoteSection{日本語訳}
\begin{itemize}
\item 試験会場と時刻は講座のWebサイトで確認できる。
\item 情報は定期的に更新される。
\item 試験日は早めに告知され、試験教室は試験直前に割り当てられる。
\item 必要な情報がWebサイトにない場合のみメールで問い合わせる。
\end{itemize}
\end{multicols}

\clearpage
\SlideHeader{Lecture 1: Overview}{010}{過去問}
\SlideImage{10}{../Materials/2026/3DM_01_Overview_SS26.pdf}
\begin{multicols}{2}\scriptsize\raggedcolumns
\NoteSection{日本語訳}
\begin{itemize}
\item 近年の過去問の印刷版は研究所で入手できる。部数には限りがある。
\item 過去問はWebサイトでも確認できる。
\item 過去問を解くだけでは十分な準備ではなく、良い成績を保証しない。
\item 講義で扱った全トピックが試験対象である。
\item 問題形式は似ることがあるが、内容は変わり、過去問にない新問も含まれ得る。
\item 過去問の模範解答は提供されない。
\end{itemize}
\end{multicols}

\clearpage
\SlideHeader{Lecture 1: Overview}{011}{アジェンダ}
\SlideImage{11}{../Materials/2026/3DM_01_Overview_SS26.pdf}
\begin{multicols}{2}\scriptsize\raggedcolumns
\NoteSection{日本語訳}
\begin{itemize}
\item 動機
\item Operations Research と Business Analytics の復習
\item モデリングと講義全体の見取り図
\end{itemize}
\end{multicols}

\clearpage
\SlideHeader{Lecture 1: Overview}{012}{動機}
\SlideImage{12}{../Materials/2026/3DM_01_Overview_SS26.pdf}
\begin{multicols}{2}\footnotesize\raggedcolumns
\NoteSection{日本語訳}
\begin{itemize}
\item 新しく変化するビジネスコンセプト: より速く、より個別的に、リアルタイムに。
\item 例
\begin{itemize}
\item 即時配送
\item オンデマンド生産・受注生産
\item 予約スケジューリング
\item その他
\end{itemize}
\item 計画と制御が難しくなる。
\item 二つの共通テーマ
\begin{itemize}
\item Stochasticity: 情報が時間とともに変化する。
\item Dynamism: 情報変化を踏まえて計画を更新する。
\end{itemize}
\item 確率的・動的意思決定過程へつながる。
\end{itemize}
\NoteSection{解説}
DDDMでは、データを眺めるだけでなく、データに基づいて実行可能な行動を選ぶ。例えば、食品配送なら『どの注文をどのドライバーに割り当てるか』、在庫管理なら『いつ何個発注するか』が決定である。\par

stochasticity は、将来情報が確実には分からないことを意味する。例として、注文発生、道路混雑、サービス時間、ドライバーのキャンセルがある。dynamism は、新しい情報が入るたびに状態と計画が変わることである。\par

試験答案では、stochasticity=不確実な情報実現、dynamism=時間を通じた状態変化と再意思決定、と分けて書き、最後に『そのため将来を先読みしつつ柔軟性を残す必要がある』と結ぶとよい。\par
\NoteSection{試験対策ポイント}
\begin{itemize}
\item DDDM = データ分析 + 意思決定モデル + 現実への実装、というループで理解する。
\item 単なる予測ではなく、予測・観測を意思決定に変換する点が講義の中心である。
\end{itemize}
\NoteSection{確認問題と解答}
\begin{enumerate}
\item \textbf{DDDMは通常のデータ分析と何が違うか。}\\通常のデータ分析は、過去データの記述や将来予測で終わる場合がある。DDDMでは、その予測や観測を使って具体的な行動を選ぶ。配送なら需要予測だけでなく、車両割当、訪問順序、注文受諾、再計画まで含める。したがって情報モデル、意思決定モデル、実装を一体で説明する必要がある。
\item \textbf{stochasticity と dynamism を例で区別せよ。}\\食品配送で『次にどこから注文が来るか分からない』ことは stochasticity である。注文が実際に入った後、ドライバー位置や未配送注文が変わり、計画を作り直す必要があることは dynamism である。両者が合わさると stochastic dynamic decision problem になる。
\end{enumerate}
\end{multicols}

\clearpage
\SlideHeader{Lecture 1: Overview}{013}{例: 輸送}
\SlideImage{13}{../Materials/2026/3DM_01_Overview_SS26.pdf}
\begin{multicols}{2}\footnotesize\raggedcolumns
\NoteSection{日本語訳}
\begin{itemize}
\item 同日配送、ピックアップ、ドローン配送、食品配送など、輸送サービスの例。
\end{itemize}
\NoteSection{解説}
DDDMでは、データを眺めるだけでなく、データに基づいて実行可能な行動を選ぶ。例えば、食品配送なら『どの注文をどのドライバーに割り当てるか』、在庫管理なら『いつ何個発注するか』が決定である。\par

stochasticity は、将来情報が確実には分からないことを意味する。例として、注文発生、道路混雑、サービス時間、ドライバーのキャンセルがある。dynamism は、新しい情報が入るたびに状態と計画が変わることである。\par

試験答案では、stochasticity=不確実な情報実現、dynamism=時間を通じた状態変化と再意思決定、と分けて書き、最後に『そのため将来を先読みしつつ柔軟性を残す必要がある』と結ぶとよい。\par
\NoteSection{試験対策ポイント}
\begin{itemize}
\item DDDM = データ分析 + 意思決定モデル + 現実への実装、というループで理解する。
\item 単なる予測ではなく、予測・観測を意思決定に変換する点が講義の中心である。
\end{itemize}
\NoteSection{確認問題と解答}
\begin{enumerate}
\item \textbf{DDDMは通常のデータ分析と何が違うか。}\\通常のデータ分析は、過去データの記述や将来予測で終わる場合がある。DDDMでは、その予測や観測を使って具体的な行動を選ぶ。配送なら需要予測だけでなく、車両割当、訪問順序、注文受諾、再計画まで含める。したがって情報モデル、意思決定モデル、実装を一体で説明する必要がある。
\item \textbf{stochasticity と dynamism を例で区別せよ。}\\食品配送で『次にどこから注文が来るか分からない』ことは stochasticity である。注文が実際に入った後、ドライバー位置や未配送注文が変わり、計画を作り直す必要があることは dynamism である。両者が合わさると stochastic dynamic decision problem になる。
\end{enumerate}
\end{multicols}

\clearpage
\SlideHeader{Lecture 1: Overview}{014}{例: モビリティとサービス}
\SlideImage{14}{../Materials/2026/3DM_01_Overview_SS26.pdf}
\begin{multicols}{2}\footnotesize\raggedcolumns
\NoteSection{日本語訳}
\begin{itemize}
\item バイクシェア、オンデマンド交通、医療訪問、タクシー、エレベーター保守、緊急車両など、モビリティ・サービスの例。
\end{itemize}
\NoteSection{解説}
DDDMでは、データを眺めるだけでなく、データに基づいて実行可能な行動を選ぶ。例えば、食品配送なら『どの注文をどのドライバーに割り当てるか』、在庫管理なら『いつ何個発注するか』が決定である。\par

stochasticity は、将来情報が確実には分からないことを意味する。例として、注文発生、道路混雑、サービス時間、ドライバーのキャンセルがある。dynamism は、新しい情報が入るたびに状態と計画が変わることである。\par

試験答案では、stochasticity=不確実な情報実現、dynamism=時間を通じた状態変化と再意思決定、と分けて書き、最後に『そのため将来を先読みしつつ柔軟性を残す必要がある』と結ぶとよい。\par
\NoteSection{試験対策ポイント}
\begin{itemize}
\item DDDM = データ分析 + 意思決定モデル + 現実への実装、というループで理解する。
\item 単なる予測ではなく、予測・観測を意思決定に変換する点が講義の中心である。
\end{itemize}
\NoteSection{確認問題と解答}
\begin{enumerate}
\item \textbf{DDDMは通常のデータ分析と何が違うか。}\\通常のデータ分析は、過去データの記述や将来予測で終わる場合がある。DDDMでは、その予測や観測を使って具体的な行動を選ぶ。配送なら需要予測だけでなく、車両割当、訪問順序、注文受諾、再計画まで含める。したがって情報モデル、意思決定モデル、実装を一体で説明する必要がある。
\item \textbf{stochasticity と dynamism を例で区別せよ。}\\食品配送で『次にどこから注文が来るか分からない』ことは stochasticity である。注文が実際に入った後、ドライバー位置や未配送注文が変わり、計画を作り直す必要があることは dynamism である。両者が合わさると stochastic dynamic decision problem になる。
\end{enumerate}
\end{multicols}

\clearpage
\SlideHeader{Lecture 1: Overview}{015}{まとめ}
\SlideImage{15}{../Materials/2026/3DM_01_Overview_SS26.pdf}
\begin{multicols}{2}\footnotesize\raggedcolumns
\NoteSection{日本語訳}
\begin{itemize}
\item サービスプロセスは時間を通じて進行する。
\item 需要、依頼、交通、ドライバー、バッテリー残量、要件などの新情報が発生する。
\item 計画への反応と更新が必要になる。
\item 反応が遅すぎる場合がある。例: 車両が渋滞に巻き込まれる、食事配送のドライバーが周辺にいない、ドローンが墜落する。
\item 期待される将来を先読みする必要がある。
\item 将来のシステム状態に対する柔軟性を目指す。
\end{itemize}
\NoteSection{解説}
DDDMでは、データを眺めるだけでなく、データに基づいて実行可能な行動を選ぶ。例えば、食品配送なら『どの注文をどのドライバーに割り当てるか』、在庫管理なら『いつ何個発注するか』が決定である。\par

stochasticity は、将来情報が確実には分からないことを意味する。例として、注文発生、道路混雑、サービス時間、ドライバーのキャンセルがある。dynamism は、新しい情報が入るたびに状態と計画が変わることである。\par

試験答案では、stochasticity=不確実な情報実現、dynamism=時間を通じた状態変化と再意思決定、と分けて書き、最後に『そのため将来を先読みしつつ柔軟性を残す必要がある』と結ぶとよい。\par
\NoteSection{試験対策ポイント}
\begin{itemize}
\item DDDM = データ分析 + 意思決定モデル + 現実への実装、というループで理解する。
\item 単なる予測ではなく、予測・観測を意思決定に変換する点が講義の中心である。
\end{itemize}
\NoteSection{確認問題と解答}
\begin{enumerate}
\item \textbf{DDDMは通常のデータ分析と何が違うか。}\\通常のデータ分析は、過去データの記述や将来予測で終わる場合がある。DDDMでは、その予測や観測を使って具体的な行動を選ぶ。配送なら需要予測だけでなく、車両割当、訪問順序、注文受諾、再計画まで含める。したがって情報モデル、意思決定モデル、実装を一体で説明する必要がある。
\item \textbf{stochasticity と dynamism を例で区別せよ。}\\食品配送で『次にどこから注文が来るか分からない』ことは stochasticity である。注文が実際に入った後、ドライバー位置や未配送注文が変わり、計画を作り直す必要があることは dynamism である。両者が合わさると stochastic dynamic decision problem になる。
\end{enumerate}
\end{multicols}

\clearpage
\SlideHeader{Lecture 1: Overview}{016}{意思決定}
\SlideImage{16}{../Materials/2026/3DM_01_Overview_SS26.pdf}
\begin{multicols}{2}\footnotesize\raggedcolumns
\NoteSection{日本語訳}
\begin{itemize}
\item Real-World: 行動が実行され、その結果として状態を観測できる。
\item Aggregation: 現実世界のデータを情報モデルの次元へ移す。
\item Information Model: 情報を圧縮された形で描写する。
\item Decision Model: 決定空間を定義し、決定を評価する。
\item Implementation: 意思決定モデルの解を現実世界の行動へ変換する。
\end{itemize}
\NoteSection{解説}
Real-World は、実際の配送現場、道路、顧客、車両、作業員のような現実対象である。現実は非常に複雑なので、そのまま最適化モデルには入れられない。\par

Aggregation は、現実データを意思決定で使える形に圧縮する処理である。例: GPS点列を道路セグメント別速度へ変換する、住所を顧客ノードへ変換する、過去走行記録から旅行時間行列を作る。\par

Information Model は『意思決定に必要な情報の表現』である。顧客集合、デポ、道路ネットワーク、旅行時間行列、需要分布、時間帯などが入る。Decision Model は『その情報を使って何を最適化するか』を表し、決定変数、目的関数、制約、実行可能解を定義する。\par
\NoteSection{試験対策ポイント}
\begin{itemize}
\item Real-World -> Aggregation -> Information Model -> Decision Model -> Implementation -> Real-World の流れを、配送例で説明できる。
\item Information Model と Decision Model を混同しない。前者は情報表現、後者は意思決定の評価・選択構造である。
\end{itemize}
\NoteSection{確認問題と解答}
\begin{enumerate}
\item \textbf{Information Model と Decision Model の違いを説明せよ。}\\Information Model は、現実を意思決定に必要な形へ圧縮した情報表現である。例として、顧客集合、距離行列、時間帯別旅行時間、需要分布がある。Decision Model は、その情報を用いて選択肢を評価する数学的・論理的構造であり、決定変数、目的関数、制約を持つ。配送例なら、旅行時間行列は情報モデル、訪問順序を選んで総旅行時間を最小化するTSP/VRPは意思決定モデルである。
\item \textbf{Aggregation と Implementation を配送例で説明せよ。}\\Aggregation は、GPS記録、道路地図、住所データなどを、顧客ノードや旅行時間行列に変換する処理である。Implementation は、最適化モデルが出した訪問順序やアーク選択を、実際のドライバーが使える配送順、住所リスト、ナビゲーションに変換する処理である。Aggregation は現実からモデルへ、Implementation はモデルから現実への橋である。
\end{enumerate}
\end{multicols}

\clearpage
\SlideHeader{Lecture 1: Overview}{017}{DDDMにおける意思決定}
\SlideImage{17}{../Materials/2026/3DM_01_Overview_SS26.pdf}
\begin{multicols}{2}\footnotesize\raggedcolumns
\NoteSection{日本語訳}
\begin{itemize}
\item DDDMでは、情報モデルだけを詳しく扱うわけではない。
\item 最適化モデルだけに焦点を当てるわけでもない。
\item 現実世界のプロセスを含むループ全体を考慮する。
\item そのため、システムを時間を通じてモデル化する必要がある。これは dynamic である。
\item 現実世界は外部効果により変化し得る。これは stochastic である。
\item 実装された決定は、すでに変化した現実世界と衝突する可能性がある。
\end{itemize}
\NoteSection{解説}
Real-World は、実際の配送現場、道路、顧客、車両、作業員のような現実対象である。現実は非常に複雑なので、そのまま最適化モデルには入れられない。\par

Aggregation は、現実データを意思決定で使える形に圧縮する処理である。例: GPS点列を道路セグメント別速度へ変換する、住所を顧客ノードへ変換する、過去走行記録から旅行時間行列を作る。\par

Information Model は『意思決定に必要な情報の表現』である。顧客集合、デポ、道路ネットワーク、旅行時間行列、需要分布、時間帯などが入る。Decision Model は『その情報を使って何を最適化するか』を表し、決定変数、目的関数、制約、実行可能解を定義する。\par
\NoteSection{試験対策ポイント}
\begin{itemize}
\item Real-World -> Aggregation -> Information Model -> Decision Model -> Implementation -> Real-World の流れを、配送例で説明できる。
\item Information Model と Decision Model を混同しない。前者は情報表現、後者は意思決定の評価・選択構造である。
\end{itemize}
\NoteSection{確認問題と解答}
\begin{enumerate}
\item \textbf{Information Model と Decision Model の違いを説明せよ。}\\Information Model は、現実を意思決定に必要な形へ圧縮した情報表現である。例として、顧客集合、距離行列、時間帯別旅行時間、需要分布がある。Decision Model は、その情報を用いて選択肢を評価する数学的・論理的構造であり、決定変数、目的関数、制約を持つ。配送例なら、旅行時間行列は情報モデル、訪問順序を選んで総旅行時間を最小化するTSP/VRPは意思決定モデルである。
\item \textbf{Aggregation と Implementation を配送例で説明せよ。}\\Aggregation は、GPS記録、道路地図、住所データなどを、顧客ノードや旅行時間行列に変換する処理である。Implementation は、最適化モデルが出した訪問順序やアーク選択を、実際のドライバーが使える配送順、住所リスト、ナビゲーションに変換する処理である。Aggregation は現実からモデルへ、Implementation はモデルから現実への橋である。
\end{enumerate}
\end{multicols}

\clearpage
\ExamHeader{Business Analytics とDDDMの全体像}
\ExamQuestion{WS23/24 Aufgabe 1 [6 Punkte]}
\ExamSubsection{問題内容}
Business Analytics が用いる三つの分野を挙げ、各二分野の交差領域も説明する問題。合計で三つの分野と三つの交差領域が問われる。\par
\ExamSubsection{満点答案}
Business Analytics は、Operations Research、Statistics、Computer Science/Information Systems の三分野を組み合わせる。Operations Research は、目的関数、制約、決定変数、実行可能解、最適化アルゴリズムを扱い、どの行動を選ぶべきかを定式化する。Statistics は、不確実なデータから分布、期待値、分散、予測、推定を扱い、観測値の背後にある規則性や不確実性を表す。Computer Science/Information Systems は、データの保存、処理、アルゴリズム実装、情報システムへの統合を扱う。\par
Operations Research と Statistics の交差では、確率的最適化や不確実性下の意思決定が中心になる。Statistics と Computer Science の交差では、データマイニング、機械学習、データ処理が中心になる。Computer Science と Operations Research の交差では、最適化アルゴリズム、計算複雑性、意思決定支援システムが中心になる。Business Analytics はこれらを統合し、データから実行可能な意思決定を導く。\par
\ExamSubsection{具体例による解説}
配送計画では、Statistics が旅行時間分布を推定し、Computer Science がGPS・注文データを処理し、Operations Research が車両ルートを最適化する。三つのうち一つでも欠けると、現実データに基づく実行可能な意思決定にならない。\par
\ExamSubsection{採点で落とさない要素}
\begin{itemize}
\item 三分野を正しく列挙
\item 三つの二分野交差を説明
\item Business Analytics が統合領域であることを明示
\item 配送などの具体例で説明
\end{itemize}
\vspace{2mm}\hrule\vspace{2mm}
\ExamQuestion{WS23/24 Aufgabe 2 [5 Punkte]}
\ExamSubsection{問題内容}
Real-World Problem と Decision Model がどのように関係し、どの概念的構成要素とデータフローで結ばれるかを説明する問題。\par
\ExamSubsection{満点答案}
現実問題は直接そのまま意思決定モデルにはならない。まず Real-World で行動が実行され、その結果として状態やデータが観測される。Aggregation は、観測データを情報モデルの次元へ変換する。Information Model は、現実を意思決定に必要な形へ圧縮した表現であり、例えば顧客、道路、旅行時間、需要などを含む。Decision Model は、この情報モデルを使い、決定空間、目的関数、制約、評価基準を定義する。Implementation または Disaggregation は、モデルの解を現実の行動へ戻す。\par
データフローは、Real-World から Aggregation を通じて Information Model へ入り、Information Model が Decision Model の入力になる。Decision Model の解は Implementation により Real-World へ戻る。実装結果を観測することで、仮説、情報モデル、意思決定モデルを更新する。したがって、関係は一方向ではなく循環である。\par
\ExamSubsection{具体例による解説}
配送では、GPSと注文データを集約して顧客集合と旅行時間行列を作る。これが情報モデルである。そこからTSP/VRPを作り、訪問順序を決める。訪問順序をドライバーのナビに変換して実行し、遅延や実走時間を再び観測してモデルを改善する。\par
\ExamSubsection{採点で落とさない要素}
\begin{itemize}
\item Real-World, Aggregation, Information Model, Decision Model, Implementationを全て説明
\item データフローを正しい順序で説明
\item ループとして説明
\item 具体例を含める
\end{itemize}
\vspace{2mm}\hrule\vspace{2mm}

\clearpage
\SlideHeader{Lecture 1: Overview}{018}{アジェンダ}
\SlideImage{18}{../Materials/2026/3DM_01_Overview_SS26.pdf}
\begin{multicols}{2}\scriptsize\raggedcolumns
\NoteSection{日本語訳}
\begin{itemize}
\item 動機
\item Operations Research と Business Analytics の復習
\item モデリングと講義全体の見取り図
\end{itemize}
\end{multicols}

\clearpage
\SlideHeader{Lecture 1: Overview}{019}{例: Braunschweigの小包配送}
\SlideImage{19}{../Materials/2026/3DM_01_Overview_SS26.pdf}
\begin{multicols}{2}\footnotesize\raggedcolumns
\NoteSection{日本語訳}
\begin{itemize}
\item トラックはデポにいる。
\item 顧客と住所がある。
\item 道路ネットワークがある。
\item 全顧客を訪問するように、市内でのドライバーのナビゲーションを探す。
\item 目標は迅速な配送である。
\item どうすれば目標を達成できるか。
\end{itemize}
\NoteSection{解説}
TSPでは、顧客をノード、顧客間移動をアーク、距離または旅行時間を重みとして表す。目的は、全顧客を一度ずつ訪問し、総移動時間または距離を最小にする訪問順序を求めることである。\par

情報モデルとしては、顧客集合、デポ、顧客間旅行時間行列が必要である。意思決定モデルとしては、アークを使うかどうかを表す二値変数 x\_\{ij\}、各顧客へ一度入る・一度出る制約、部分巡回除去制約、総旅行時間最小化の目的関数が必要である。\par

時間依存を入れると、旅行時間は定数 d\_\{ij\} ではなく、出発時刻 t に依存する d\_\{ij\}(t) になる。ある顧客を先に訪問するか後に訪問するかで、後続アークの出発時刻が変わり、旅行時間も変わる。\par
\NoteSection{試験対策ポイント}
\begin{itemize}
\item TSPの基本定式化を、変数・目的関数・制約に分けて説明できる。
\item 時間依存TSPでは、旅行時間行列だけでなく時刻依存関数や複数行列が必要になる。
\end{itemize}
\NoteSection{確認問題と解答}
\begin{enumerate}
\item \textbf{TSPの数理モデルを言葉で説明せよ。}\\ノード集合はデポと顧客、アーク集合は顧客間移動を表す。二値変数 x\_\{ij\}=1 は、i から j へ直接移動することを意味する。目的は sum d\_\{ij\}x\_\{ij\} を最小化することである。制約として、各顧客に一度入る、各顧客から一度出る、複数の小さな巡回ができないようにする部分巡回除去制約が必要である。
\item \textbf{時間依存旅行時間がTSPを難しくする理由を説明せよ。}\\静的TSPではアークのコストは固定である。しかし時間依存TSPでは、iからjへのコストが出発時刻に依存する。ある挿入や訪問順変更により後続の到着時刻が変わり、その後の全アークの旅行時間も変わるため、局所的な評価が全体に波及する。
\end{enumerate}
\end{multicols}

\clearpage
\SlideHeader{Lecture 1: Overview}{020}{モデリング: 巡回セールスマン問題}
\SlideImage{20}{../Materials/2026/3DM_01_Overview_SS26.pdf}
\begin{multicols}{2}\footnotesize\raggedcolumns
\NoteSection{日本語訳}
\begin{itemize}
\item Part 1
\begin{itemize}
\item 道路ネットワークを顧客間のアークへ変換する。
\item 走行速度をアークごとの旅行時間へ変換する。
\item 重み付きグラフまたは旅行時間行列を得る。
\end{itemize}
\item Part 2
\begin{itemize}
\item ナビゲーションを、使用アークに関する決定変数へ変換する。
\item 実行可能なナビゲーションを保証する制約を追加する。
\item Part 1に基づいて目的関数を定義する。
\end{itemize}
\end{itemize}
\NoteSection{解説}
TSPでは、顧客をノード、顧客間移動をアーク、距離または旅行時間を重みとして表す。目的は、全顧客を一度ずつ訪問し、総移動時間または距離を最小にする訪問順序を求めることである。\par

情報モデルとしては、顧客集合、デポ、顧客間旅行時間行列が必要である。意思決定モデルとしては、アークを使うかどうかを表す二値変数 x\_\{ij\}、各顧客へ一度入る・一度出る制約、部分巡回除去制約、総旅行時間最小化の目的関数が必要である。\par

時間依存を入れると、旅行時間は定数 d\_\{ij\} ではなく、出発時刻 t に依存する d\_\{ij\}(t) になる。ある顧客を先に訪問するか後に訪問するかで、後続アークの出発時刻が変わり、旅行時間も変わる。\par
\NoteSection{試験対策ポイント}
\begin{itemize}
\item TSPの基本定式化を、変数・目的関数・制約に分けて説明できる。
\item 時間依存TSPでは、旅行時間行列だけでなく時刻依存関数や複数行列が必要になる。
\end{itemize}
\NoteSection{確認問題と解答}
\begin{enumerate}
\item \textbf{TSPの数理モデルを言葉で説明せよ。}\\ノード集合はデポと顧客、アーク集合は顧客間移動を表す。二値変数 x\_\{ij\}=1 は、i から j へ直接移動することを意味する。目的は sum d\_\{ij\}x\_\{ij\} を最小化することである。制約として、各顧客に一度入る、各顧客から一度出る、複数の小さな巡回ができないようにする部分巡回除去制約が必要である。
\item \textbf{時間依存旅行時間がTSPを難しくする理由を説明せよ。}\\静的TSPではアークのコストは固定である。しかし時間依存TSPでは、iからjへのコストが出発時刻に依存する。ある挿入や訪問順変更により後続の到着時刻が変わり、その後の全アークの旅行時間も変わるため、局所的な評価が全体に波及する。
\end{enumerate}
\end{multicols}

\clearpage
\SlideHeader{Lecture 1: Overview}{021}{解を見つける}
\SlideImage{21}{../Materials/2026/3DM_01_Overview_SS26.pdf}
\begin{multicols}{2}\footnotesize\raggedcolumns
\NoteSection{日本語訳}
\begin{itemize}
\item モデルは解空間を張り、解を評価する。
\item 候補解の求め方
\begin{itemize}
\item グラフベース手法、例: insertion procedures
\item 混合整数計画、例: branch and bound
\end{itemize}
\item 最後のステップ: 実務への実装
\begin{itemize}
\item 数字を顧客へ変換する。
\item 選択されたアークを道路セグメントの集合へ対応させる。
\item ドライバーにナビゲーションを提供する。
\end{itemize}
\end{itemize}
\NoteSection{解説}
TSPでは、顧客をノード、顧客間移動をアーク、距離または旅行時間を重みとして表す。目的は、全顧客を一度ずつ訪問し、総移動時間または距離を最小にする訪問順序を求めることである。\par

情報モデルとしては、顧客集合、デポ、顧客間旅行時間行列が必要である。意思決定モデルとしては、アークを使うかどうかを表す二値変数 x\_\{ij\}、各顧客へ一度入る・一度出る制約、部分巡回除去制約、総旅行時間最小化の目的関数が必要である。\par

時間依存を入れると、旅行時間は定数 d\_\{ij\} ではなく、出発時刻 t に依存する d\_\{ij\}(t) になる。ある顧客を先に訪問するか後に訪問するかで、後続アークの出発時刻が変わり、旅行時間も変わる。\par
\NoteSection{試験対策ポイント}
\begin{itemize}
\item TSPの基本定式化を、変数・目的関数・制約に分けて説明できる。
\item 時間依存TSPでは、旅行時間行列だけでなく時刻依存関数や複数行列が必要になる。
\end{itemize}
\NoteSection{確認問題と解答}
\begin{enumerate}
\item \textbf{TSPの数理モデルを言葉で説明せよ。}\\ノード集合はデポと顧客、アーク集合は顧客間移動を表す。二値変数 x\_\{ij\}=1 は、i から j へ直接移動することを意味する。目的は sum d\_\{ij\}x\_\{ij\} を最小化することである。制約として、各顧客に一度入る、各顧客から一度出る、複数の小さな巡回ができないようにする部分巡回除去制約が必要である。
\item \textbf{時間依存旅行時間がTSPを難しくする理由を説明せよ。}\\静的TSPではアークのコストは固定である。しかし時間依存TSPでは、iからjへのコストが出発時刻に依存する。ある挿入や訪問順変更により後続の到着時刻が変わり、その後の全アークの旅行時間も変わるため、局所的な評価が全体に波及する。
\end{enumerate}
\end{multicols}

\clearpage
\SlideHeader{Lecture 1: Overview}{022}{Operations Researchの三段階}
\SlideImage{22}{../Materials/2026/3DM_01_Overview_SS26.pdf}
\begin{multicols}{2}\footnotesize\raggedcolumns
\NoteSection{日本語訳}
\begin{itemize}
\item 問題定義
\item モデリング
\item 解法
\item 情報と意思決定
\item 実装を含む。
\end{itemize}
\NoteSection{解説}
IDAは appearance から structure へ向かう。つまり、記録データの表面的なパターンから、背後の情報構造を推測する。例: GPS走行ログから道路混雑パターンを推定する。\par

PMT/ORは structure から appearance/solution へ向かう。つまり、問題構造を仮定し、目的関数・制約・決定変数を作り、解を求める。例: 顧客集合と旅行時間行列からTSPを作る。\par

Business Analytics では、データ側と意思決定側をつなぐ。データ分析だけでは『何を実行すべきか』が出ない。最適化だけでも、現実のデータを正しく表現しなければ役に立たない。\par
\NoteSection{試験対策ポイント}
\begin{itemize}
\item IDA = データから構造を発見する方向、PMT/OR = 構造から意思決定モデルと解を作る方向。
\item 三分野を列挙するだけでなく、交差領域がなぜ必要かを説明する。
\end{itemize}
\NoteSection{確認問題と解答}
\begin{enumerate}
\item \textbf{IDA と PMT/OR の方向性の違いを説明せよ。}\\IDAは記録データ、つまり appearance から出発し、クリーニング、補完、検証、データマイニングを通じて情報構造やパターンを見つける。PMT/ORは問題構造を仮説として置き、目的、制約、決定空間を定義して、実行可能な解を計算する。DDDMでは両方が必要で、情報モデルが両者を接続する。
\item \textbf{Business Analytics の三分野を使って、配送計画を説明せよ。}\\Statistics は旅行時間や需要の不確実性を推定する。Computer Science/Information Systems はGPS、注文、地図データを保存・処理し、システムとして実装する。Operations Research は訪問順序、車両割当、制約を最適化する。三つを統合することで、データに基づく実行可能な配送計画が得られる。
\end{enumerate}
\end{multicols}

\clearpage
\SlideHeader{Lecture 1: Overview}{023}{アジェンダ}
\SlideImage{23}{../Materials/2026/3DM_01_Overview_SS26.pdf}
\begin{multicols}{2}\scriptsize\raggedcolumns
\NoteSection{日本語訳}
\begin{itemize}
\item 動機
\item Operations Research と Business Analytics の復習
\item モデリングと講義全体の見取り図
\end{itemize}
\end{multicols}

\clearpage
\SlideHeader{Lecture 1: Overview}{024}{モデリング}
\SlideImage{24}{../Materials/2026/3DM_01_Overview_SS26.pdf}
\begin{multicols}{2}\footnotesize\raggedcolumns
\NoteSection{日本語訳}
\begin{itemize}
\item 異なる問題が同じモデルを共有する場合がある。例: 小包配送と技術者サービス。
\item 同じ問題を異なる方法でモデル化することもできる。
\item 情報モデルと意思決定モデルの間には接続がある。
\end{itemize}
\NoteSection{解説}
Real-World は、実際の配送現場、道路、顧客、車両、作業員のような現実対象である。現実は非常に複雑なので、そのまま最適化モデルには入れられない。\par

Aggregation は、現実データを意思決定で使える形に圧縮する処理である。例: GPS点列を道路セグメント別速度へ変換する、住所を顧客ノードへ変換する、過去走行記録から旅行時間行列を作る。\par

Information Model は『意思決定に必要な情報の表現』である。顧客集合、デポ、道路ネットワーク、旅行時間行列、需要分布、時間帯などが入る。Decision Model は『その情報を使って何を最適化するか』を表し、決定変数、目的関数、制約、実行可能解を定義する。\par
\NoteSection{試験対策ポイント}
\begin{itemize}
\item Real-World -> Aggregation -> Information Model -> Decision Model -> Implementation -> Real-World の流れを、配送例で説明できる。
\item Information Model と Decision Model を混同しない。前者は情報表現、後者は意思決定の評価・選択構造である。
\end{itemize}
\NoteSection{確認問題と解答}
\begin{enumerate}
\item \textbf{Information Model と Decision Model の違いを説明せよ。}\\Information Model は、現実を意思決定に必要な形へ圧縮した情報表現である。例として、顧客集合、距離行列、時間帯別旅行時間、需要分布がある。Decision Model は、その情報を用いて選択肢を評価する数学的・論理的構造であり、決定変数、目的関数、制約を持つ。配送例なら、旅行時間行列は情報モデル、訪問順序を選んで総旅行時間を最小化するTSP/VRPは意思決定モデルである。
\item \textbf{Aggregation と Implementation を配送例で説明せよ。}\\Aggregation は、GPS記録、道路地図、住所データなどを、顧客ノードや旅行時間行列に変換する処理である。Implementation は、最適化モデルが出した訪問順序やアーク選択を、実際のドライバーが使える配送順、住所リスト、ナビゲーションに変換する処理である。Aggregation は現実からモデルへ、Implementation はモデルから現実への橋である。
\end{enumerate}
\end{multicols}

\clearpage
\SlideHeader{Lecture 1: Overview}{025}{例I: 代替的な情報モデル}
\SlideImage{25}{../Materials/2026/3DM_01_Overview_SS26.pdf}
\begin{multicols}{2}\footnotesize\raggedcolumns
\NoteSection{日本語訳}
\begin{itemize}
\item 旅行時間の日内変動、つまり時間依存旅行時間を観測する。
\item 7-9時、9-15時、15-18時のような時間帯を考える。
\item 情報モデルは複数の旅行時間行列の集合を含む。
\item 情報生成はより複雑になる。
\item 代替的な意思決定モデルが必要になる。
\end{itemize}
\NoteSection{解説}
時間依存旅行時間では、道路や顧客間の旅行時間が時刻で変わる。朝の通勤時間、昼間、夕方では、同じ道路でも速度が違う。したがって一つの平均旅行時間だけでは現実を十分に表せない。\par

情報モデルとしては、時間帯別旅行時間行列、または連続的な旅行時間関数を使う。例: 7-9時、9-15時、15-18時で別の行列を持つ。\par

重要なのは、時間依存性は情報モデルだけでなく意思決定モデルも変える点である。訪問順序を決めるだけでなく、どの時刻にどのアークを通るかが決定に入る。\par
\NoteSection{試験対策ポイント}
\begin{itemize}
\item 詳細な情報モデルほど現実に近いが、データ量・計算複雑性・解釈可能性にコストがある。
\item 時間依存旅行時間は、FCDなどから推定される。
\end{itemize}
\NoteSection{確認問題と解答}
\begin{enumerate}
\item \textbf{なぜ常に最も詳細な旅行時間モデルを選ばないのか。}\\詳細なモデルは時間帯・道路セグメントごとに多くの観測を必要とする。観測が少ないセルでは平均値が不安定になり、過学習のように現実を誤って表す可能性がある。また、意思決定モデルの変数や制約が増え、計算時間が長くなる。したがって、意思決定に有用な粒度を選ぶ必要がある。
\item \textbf{時間帯別旅行時間行列の長所と短所を述べよ。}\\長所は、朝夕の混雑など平均行列では消える変動を反映できることである。短所は、時間帯ごとに十分な観測が必要であり、境界時刻で旅行時間が不連続になる可能性があること、意思決定モデルが複雑になることである。
\end{enumerate}
\end{multicols}

\clearpage
\SlideHeader{Lecture 1: Overview}{026}{例II: 代替的な意思決定モデル}
\SlideImage{26}{../Materials/2026/3DM_01_Overview_SS26.pdf}
\begin{multicols}{2}\footnotesize\raggedcolumns
\NoteSection{日本語訳}
\begin{itemize}
\item 解が実行可能に見えても、残業が頻繁に発生する。
\item 旅行時間に不確実性がある。
\item 意思決定モデルを変更する。これは stochasticity に関係する。
\item 代替的でより詳細な情報モデルが必要になる。
\end{itemize}
\NoteSection{解説}
不確実性とは、意思決定時点では将来の実現値が完全には分からないことである。需要、交通、サービス時間、ドライバーの可用性、天候、政治状況などが例である。\par

不確実性は、頻度、範囲、破壊力、予測可能性で異なる。例えば毎日変わる交通量は高頻度だが比較的予測可能であり、突然の道路閉鎖は低頻度だが破壊力が大きい。\par

不確実性は情報モデルに入るだけでなく、意思決定モデルにも影響する。平均値だけで計画すると、遅延リスクやサービス品質低下を見落とす可能性がある。\par
\NoteSection{試験対策ポイント}
\begin{itemize}
\item 不確実性の発生源を、需要・リソース・環境に分けて説明できる。
\end{itemize}
\NoteSection{確認問題と解答}
\begin{enumerate}
\item \textbf{Same-Day配送における三種類の不確実性を例示せよ。}\\需要不確実性は、注文がいつ・どこで・どれだけ発生するかである。リソース不確実性は、ドライバーの出勤、車両故障、バッテリー残量、スキルである。環境不確実性は、交通、駐車、天候、道路工事、ネットワーク障害である。
\item \textbf{平均値だけを使う計画の危険性を説明せよ。}\\平均旅行時間では、遅延の裾の長さや極端な渋滞を見落とす。全アークで平均を使うと、各区間では小さな誤差でも全体ツアーでは大きな遅延になることがある。サービス時間や交通が右裾の長い分布を持つ場合、平均値だけでは信頼性を評価できない。
\end{enumerate}
\end{multicols}

\clearpage
\SlideHeader{Lecture 1: Overview}{027}{例III: 動的性}
\SlideImage{27}{../Materials/2026/3DM_01_Overview_SS26.pdf}
\begin{multicols}{2}\footnotesize\raggedcolumns
\NoteSection{日本語訳}
\begin{itemize}
\item 情報モデルが時間とともに変化することを観測する。
\item 変化に反応できる。
\item 確率的・動的モデルが必要になる。
\item 将来の先読みが重要になる。
\item Approximate Dynamic Programming へつながる。
\end{itemize}
\NoteSection{解説}
DDDMでは、データを眺めるだけでなく、データに基づいて実行可能な行動を選ぶ。例えば、食品配送なら『どの注文をどのドライバーに割り当てるか』、在庫管理なら『いつ何個発注するか』が決定である。\par

stochasticity は、将来情報が確実には分からないことを意味する。例として、注文発生、道路混雑、サービス時間、ドライバーのキャンセルがある。dynamism は、新しい情報が入るたびに状態と計画が変わることである。\par

試験答案では、stochasticity=不確実な情報実現、dynamism=時間を通じた状態変化と再意思決定、と分けて書き、最後に『そのため将来を先読みしつつ柔軟性を残す必要がある』と結ぶとよい。\par
\NoteSection{試験対策ポイント}
\begin{itemize}
\item DDDM = データ分析 + 意思決定モデル + 現実への実装、というループで理解する。
\item 単なる予測ではなく、予測・観測を意思決定に変換する点が講義の中心である。
\end{itemize}
\NoteSection{確認問題と解答}
\begin{enumerate}
\item \textbf{DDDMは通常のデータ分析と何が違うか。}\\通常のデータ分析は、過去データの記述や将来予測で終わる場合がある。DDDMでは、その予測や観測を使って具体的な行動を選ぶ。配送なら需要予測だけでなく、車両割当、訪問順序、注文受諾、再計画まで含める。したがって情報モデル、意思決定モデル、実装を一体で説明する必要がある。
\item \textbf{stochasticity と dynamism を例で区別せよ。}\\食品配送で『次にどこから注文が来るか分からない』ことは stochasticity である。注文が実際に入った後、ドライバー位置や未配送注文が変わり、計画を作り直す必要があることは dynamism である。両者が合わさると stochastic dynamic decision problem になる。
\end{enumerate}
\end{multicols}

\clearpage
\SlideHeader{Lecture 1: Overview}{028}{次回予告}
\SlideImage{28}{../Materials/2026/3DM_01_Overview_SS26.pdf}
\begin{multicols}{2}\scriptsize\raggedcolumns
\NoteSection{日本語訳}
\begin{itemize}
\item Business Analytics
\item 情報モデリングと意思決定モデリング
\item ケーススタディ: 時間依存旅行時間を持つ配送ルーティング
\end{itemize}
\end{multicols}
